\section{Further Literature Review}
\label{sec:literature}
The PPR reviewed the origins of LMS and adaptive equalization, LMS tracking theory, decision-directed and normalised adaptation, and FPGA and fixed-point design. For the mid-term the literature was organised by development phase, so that each phase rests on its sources and each source is tied to the result it supports or contradicts (\cref{tab:litphase}). References already used in the PPR, and what they establish, are printed in black; \nb{sources added since the PPR, and every statement about our own results, are in blue. New sources were read at least at catalogue or abstract level; the depth is stated where it matters, and ``supports'' means that the source states the method or result we use, not that its derivations were re-checked.}

\noindent\textbf{P1 --- static LMS and validation.} The algorithm \cite{widrowhoff1960}, adaptive equalization \cite{lucky1965,qureshi1985}, the Wiener solution, misadjustment, eigenvalue spread and persistent excitation \cite{widrowstearns1985,haykin2014}, the BPSK error rate and the ISI model \cite{proakis2008}, and the equalizer delay set by the reference tap \cite{qureshi1985}. \nb{Mazo \cite{mazo1979} showed that the independence assumption on which LMS convergence theory rests is ``not even approximately true'' for a tapped-delay-line equalizer, although it often predicts well (abstract level). Our misadjustment finding (\cref{sec:misadj}) is a case where a consequence of that assumption --- a white residual error --- fails visibly, by a factor of 1.47. The PRBS and the noise LFSRs use maximal-length taps from Alfke \cite{alfke1996}.}

\noindent\textbf{P2 --- abrupt transitions.} Widrow \etal \cite{widrow1976} and Haykin \cite{haykin2014} give the learning-curve time constants $1/(\mu\lambda_i)$ and the role of the eigenvalue spread. \nb{Our recovery times follow $N_{\rm rec}\propto1/\mu$, and the slowest mode ($\lambda_{\min}\approx0.07$) explains why abrupt-transition runs need a longer training (\cref{sec:p2}).}

\noindent\textbf{P3--P4 --- continuous variation and the optimum step.} The gradient-noise/lag decomposition for a first-order Markov (random-walk) optimum \cite{widrow1976}, tracking in transmission \cite{macchi1995}, RLS tracking compared with LMS \cite{eleftheriou1986}, and statistical channel models \cite{jakes1974}. \nb{(The PPR described \cite{eleftheriou1986} as comparing ``across algorithm families''; it analyses RLS tracking and compares it with LMS.) For deterministic variation, Vlaski, Rizk and Sayed \cite{vlaski2020} show that a non-zero-mean drift turns the $1/\mu$ lag term into $1/\mu^2$ (read in full); Bershad and Macchi \cite{bershad1989} and an associated report \cite{dticada208618} treat LMS tracking of a chirp, with lag of order $v^{-2}$ (abstract level); Yousef and Sayed \cite{yousef2001} unify steady-state and tracking analysis; Bergmans \cite{bergmans1990} derives the lag error of LMS for a channel with band-limited gain variations and shows that the total error either rises monotonically with $\mu$ or has an interior minimum, depending on the size and bandwidth of the variation (read in full; replicated as R7). These sources give the 2/3 ramp alternative and the basis of the model of \cref{sec:theory}, which explains our measured exponent of 0.52--0.59. The statistics use the bootstrap \cite{efron1993} and the Wilson interval \cite{wilson1927}.}

\noindent\textbf{Adaptation source.} Sato \cite{sato1975} and Godard \cite{godard1980} established adaptation without a continuous training reference. \nb{Mazo \cite{mazo1980} analyses decision-directed start-up and notes that, because the symbol estimates can be wrong, convergence is not guaranteed (abstract level). Our measured divergence map (\cref{sec:dd}) quantifies this for our channel: divergence is set by $\mu$ and SNR, not by the variation rate.}

\noindent\textbf{P5 --- payloads.} \nb{The audio format (8\,kHz, 8 bits per sample) is the telephone-speech rate of ITU-T G.711 \cite{g711}; our samples are linear, not companded. The test image is a public-domain photograph distributed with scikit-image \cite{vanderwalt2014}.}

\noindent\textbf{P6 --- fixed point.} Delayed adaptation and its step bound \cite{long1989}, coefficient drift from bias in a digital update \cite{gitlin1982}, word-length selection \cite{constantinides2003}, FPGA LMS architectures \cite{hanson2016} and the DSP48E1 slice \cite{xilinxug479}. \nb{Ghanassi, Champagne and Kabal \cite{ghanassi2007} derive the steady-state error of fixed-point LMS for a stationary target, including a stall mode whose error grows as $1/\mu^2$; Shahdoosti \cite{shahdoosti2018} states the stall rule $|\mu e u|<\mathrm{LSB}/2$; Montenegro \etal \cite{montenegro2006} model stalling with Markov chains (record only); Khan and Gustafsson \cite{khan2023} study power-of-two step approximations. Our word-length study (\cref{sec:fxp}) is the tracking counterpart: the accumulator width, not the multiplier width, sets the boundary, and truncation needs about two more bits --- the drift that \cite{gitlin1982} anticipates. The sum of uniforms of the noise generator follows the Irwin--Hall distribution \cite{irwin1927,hall1927}, whose bounded tails we measured; Lee \etal \cite{lee2006} give the established higher-accuracy alternative. The FPGA LMS equalizers found were tested on static channels \cite{vargas2007} or used other algorithms \cite{kumari2021}; one thesis combines fixed-point LMS with time-varying channels on a DSP processor rather than an FPGA \cite{rierapalou2002}.}

\subsection*{Three key sources in detail}
\label{sec:keysources}
Three sources carry most of the weight of this report. For each, what it establishes and assumes is stated first, then \nb{what was reproduced and how our work goes beyond it. The type of agreement is named explicitly: \emph{exact numerical replication} (a published number is matched), \emph{reproduction of a theoretical trend} (a published functional form or scaling is matched), \emph{partial replication}, and \emph{independent validation} (our own model confirmed by another tool).}

\noindent\textbf{Widrow \etal \cite{widrow1976}} establish the learning behaviour of LMS for stationary and nonstationary inputs: the misadjustment due to gradient noise, and, for an optimum that drifts, the split of the excess error into gradient noise rising with $\mu$ and lag falling with $\mu$, with an optimum step between them, \eqref{eq:mis}. \nb{(Read in full.) The paper writes the update as $\mathbf w+2\mu e\mathbf u$, so its $\mu$ is half of ours; its results then read: misadjustment $\mu\,\mathrm{tr}(\mathbf R)$ (its Eq.~51, our $\mu\,\mathrm{tr}(\mathbf R)/2$), lag misadjustment $n\sigma^2/(4\xi_{\min}\mu)$ (Eq.~81), and $\mu^\star=[n\sigma^2/(4\xi_{\min}\mathrm{tr}\mathbf R)]^{1/2}$ (Eq.~86), which is exactly \eqref{eq:mis} with $\mathrm{tr}(\mathbf Q)=n\sigma^2$. Its assumptions are stated explicitly: input vectors uncorrelated over time and error independent of the input (Sec.~VIII); the formula valid for misadjustment up to about 25\,\%; a stationary input ($\mathbf R$ constant); each optimal weight an independent, slowly varying first-order Markov process, with the nonstationarity much slower than the adaptation (Sec.~XI); and white input, $\mathbf R=\mathbf I$, four weights and a 128-member ensemble in its tracking experiments. Two passages bear directly on our results. In its adaptive-predictor experiment the measured misadjustment exceeds the formula by 1.2--1.9 times (its Table~I; e.g.\ 3\,\% against 2.5\,\%), which the authors attribute to correlated input samples violating the uncorrelated-input assumption; our correction \eqref{eq:corrmis} quantifies that effect for an equalizer (\cref{sec:misadj}). Its footnote~4 notes that for a weight fluctuation with a different (uniform low-pass) spectrum the lag term changes its dependence on $\mu$, the optimum moves to where the gradient-noise loss is twice the lag loss, and the lag is ``quite sensitive to the spectral characteristics of the fluctuation'' --- the dependence we measure as a channel- and trajectory-dependent exponent (\cref{sec:theory}). Reproduced as theoretical trends: convergence to the Wiener filter, the misadjustment on a white input within 1--3\,\%, and the random-walk trade-off with $\mu_{\rm opt}$ within 1--9\,\% and exponent 1.000 (R3, R6). Beyond it: an ISI channel, where the white-residual assumption fails by a factor 1.47, and a deterministic, reversing ramp instead of a Markov drift.}

\noindent\nb{\textbf{Vlaski, Rizk and Sayed \cite{vlaski2020}} (read in full) analyse the steady-state tracking of stochastic-gradient learners, LMS included, when the optimum drifts. With a zero-mean (random-walk) drift the lag term scales as $1/\mu$; with a non-zero-mean, deterministic drift it scales as $1/\mu^2$; combined with the gradient-noise term, this gives an optimum-step exponent of 2/3 for a steady ramp. The analysis is first-order in a small step size and describes the steady state of a persistent drift. Reproduced as a theoretical trend: the ramp lag slope $-2.01$ and, with the exact gradient-noise term, the ramp optimum within 1--2.4\,\% (R6). Beyond it: our triangle reverses before the slowest eigen-direction of $\mathbf R$ reaches steady state, so the lag falls more slowly than $1/\mu^2$ and the measured exponent lies below 2/3 (0.52 primary, 0.58--0.59 control; \cref{sec:theory}).}

\noindent\nb{\textbf{Ghanassi, Champagne and Kabal \cite{ghanassi2007}} (read in full) derive the steady-state mean squared error of fixed-point LMS over the whole step-size range by splitting it into three modes: a stall mode, in which the quantised update rounds to zero and the error falls as $1/\mu^2$ (with two constants fitted to their simulations); a new intermediate mode; and the stochastic-gradient mode, in which quantisation acts as additive white noise. They assume a white stationary Gaussian input, a long filter, round-off quantisation of every product with no overflow, and a stationary system-identification target; their experiments use $N=100$, 16 bits and no additive noise. Reproduced: after correcting our rebuild to their two-quantiser update, our simulation matches their own simulation points within 0.72--1.19 for $\mu\le0.155$, except at the narrow error minimum with a wide accumulator (R5). Their formulas match their own points, and ours, only in parts (stall law 0.70--0.97, other modes within about a factor of two), so that difference belongs to the analysis, not to the replication. Beyond it: their analysis has no moving channel. Our word-length study measures how the same stalling mechanism moves the \emph{tracking} boundary of an equalizer on a time-varying ISI channel, and separates accumulator width from multiplier width and rounding from truncation (\cref{sec:fxp}).}

\noindent\nb{The remaining agreements are of the other two types: exact numerical replication for the BPSK BER and the stall onset at iteration 134 (R1, R4), and independent validation of our own Python LMS model, whose MSE and final coefficients a plain MATLAB implementation reproduces to every printed digit and MathWorks' equalizer within 1.2\,\% (\cref{sec:replication}).}

\noindent\textbf{Position.} The PPR observed that tracking analyses usually assume infinite precision and FPGA studies usually assume a static channel, and placed this project at their intersection. \nb{The further search did not find a study combining a bit-true fixed-point LMS equalizer, a swept deterministic variation rate, BER and MSE together, and a tracking limit. The search was short and not systematic, so this is stated only to the best of our knowledge; ``first'' or ``novel'' is not claimed.}

\noindent\nb{\textbf{Reference audit.} The audit checks that every claim the PPR makes from the literature rests on a source that exists and says what the PPR attributes to it. All 22 PPR references were checked against catalogue, publisher or author records (bibliographic details, and content at abstract level only): 10 were fully confirmed, 2 confirmed through a secondary citation, 7 partly confirmed and 3 not confirmed online; none was found wrong. Two papers were afterwards read in full, Widrow \etal \cite{widrow1976} (above) and Qureshi \cite{qureshi1985}, which is now confirmed (\emph{Proc.\ IEEE} 73(9), pp.~1349--1387). The audit led to two corrections: Qureshi does not single out non-minimum-phase channels, as the PPR stated, so that point now rests on our own delay sweep (\cref{sec:p1}); and the UG479 entry now carries its version and date. Nagumo--Noda \cite{nagumo1967} and the board manual \cite{realdigitalmanual} still need a library check.}

\begin{table*}[t]
\centering\footnotesize
\setlength{\tabcolsep}{3pt}
\caption{Literature by development phase. Black: references of the PPR and what they establish; \nb{blue: references added since the PPR, and our results.}}
\label{tab:litphase}
\begin{tabular}{@{}p{0.11\linewidth}>{\leavevmode}p{0.25\linewidth}>{\leavevmode}p{0.29\linewidth}>{\leavevmode\color{blue}}p{0.3\linewidth}@{}}
\toprule
Phase & References & What they establish & Our result (section)\\
\midrule
P1 & \cite{widrowhoff1960,lucky1965,qureshi1985,widrowstearns1985,haykin2014,proakis2008} \nb{\cite{mazo1979,alfke1996}} & LMS, Wiener solution, misadjustment, eigenvalue spread, BPSK BER, delay for non-minimum-phase channels; \nb{independence assumption; LFSR taps} & BER on theory; Wiener convergence; misadjustment 1.47$\times$ textbook (\cref{sec:p1})\\
P2 & \cite{widrow1976,haykin2014} & Time constants $1/(\mu\lambda_i)$ & $N_{\rm rec}\propto1/\mu$; slow-mode training (\cref{sec:p2})\\
P3--P4 & \cite{widrow1976,macchi1995,eleftheriou1986,jakes1974} \nb{\cite{vlaski2020,bershad1989,dticada208618,yousef2001,bergmans1990,efron1993,wilson1927}} & Lag vs gradient noise, optimum step; \nb{deterministic drift ($1/\mu^2$); bootstrap and binomial intervals} & Boundary $\approx0.11/\mu$; exponent 0.52--0.59, explained (\cref{sec:p3,sec:theory})\\
Adaptation & \cite{sato1975,godard1980} \nb{\cite{mazo1980}} & Adaptation without training; \nb{decision-directed convergence not guaranteed} & Divergence set by $\mu$ and SNR (\cref{sec:dd})\\
P5 & \nb{\cite{g711,vanderwalt2014}} & \nb{Telephone speech format; test image} & Audio and image results (\cref{sec:audio,sec:image})\\
P6 & \cite{long1989,gitlin1982,constantinides2003,hanson2016,xilinxug479} \nb{\cite{ghanassi2007,shahdoosti2018,montenegro2006,khan2023,irwin1927,hall1927,lee2006,vargas2007}} & Delayed LMS, coefficient drift, word length, FPGA LMS, DSP slice; \nb{stall mode and stall rule; sum-of-uniform noise; accurate hardware Gaussian noise} & Accumulator width sets the boundary; truncation needs $\approx$2 more bits; Irwin--Hall BER (\cref{sec:fxp,sec:replication})\\
\bottomrule
\end{tabular}
\end{table*}
